\subsection{Amdahl's law and parallel overhead (RQ6)}
We fitted $S(p)=1/(f+(1-f)/p)$ to the measured speedups of each batch size (Figure~\ref{fig:amdahl}). On the Xeon the fitted serial fraction is only $f=$ \val{Sp:amdahl:f:64} for batch 64 and \val{Sp:amdahl:f:4096} for batch 4096. This would allow at most \val{Sp:amdahl:ceil:64}$\times$ and \val{Sp:amdahl:ceil:4096}$\times$, so the measured \val{Sp:omp:speedup:4096:24}$\times$ at 24 threads is far from the Amdahl ceiling and the serial fraction is not what limits OpenMP here. The Karp--Flatt metric agrees: at batch 4096 it falls from \val{Sp:kf:4096:2} with 2 threads to \val{Sp:kf:4096:24} with 24 threads, so there is no growing overhead term. For one image per batch the fit gives $f=$ \val{Sp:amdahl:f:1}, a ceiling of only \val{Sp:amdahl:ceil:1}$\times$, and the Karp--Flatt metric stays between \val{Sp:kf:1:8} and \val{Sp:kf:1:12}: shrinking the batch increases the effective serial fraction, as Gustafson~\cite{gustafson1988reevaluating} describes for the reverse case of growing the problem. On the M4 the Karp--Flatt metric at batch 4096 jumps from \val{Mac:kf:4096:4} with 4 threads to \val{Mac:kf:4096:6} with 6 threads, where the efficiency cores join in. A heterogeneous chip violates the equal-processor assumption behind Amdahl's law, so the M4 value $f=$ \val{Mac:amdahl:f:4096} (R$^2=$ \val{Mac:amdahl:r2:4096}) is only a summary of the curve.

\widefigure{amdahl_combined.pdf}{Measured speedups (dots) and fitted Amdahl curves for four batch sizes on the Apple M4 (left) and the Xeon (right); $f$ is the fitted serial fraction.}{fig:amdahl}

Synchronisation explains why OpenMP loses for single images on the M4 and wins on the Xeon. One image needs \val{Sp:serial:usimg:1}\,$\mu$s of serial work on the Xeon but only \val{Mac:serial:usimg:1}\,$\mu$s on the M4, while one OpenMP fork--join costs \val{Sp:micro:openmp_fork_join:24}\,$\mu$s with 24 threads on the Xeon and \val{Mac:micro:openmp_fork_join:10}\,$\mu$s with 10 threads on the M4 (Table~\ref{tab:micro}). When the work per batch is comparable to the cost of synchronising the threads, as on the M4, parallelisation loses; when it is much larger, as on the Xeon, it wins even for one image. A larger batch spreads the synchronisation cost over more images, which is why OpenMP improves with the batch size.

\subsection{Memory traffic and the GPU}
If every image had to read the 803{,}840-byte hidden-layer weights from device memory, one pass would move \val{X:weights:GB}\,GB. At the \val{Sp:micro:d2d_copy_bandwidth:0}\,GB/s measured with a device-to-device copy this would take about 42\,ms, longer than the measured kernel time of \val{Sp:cuda:kernel:4096}\,ms at batch 4096; the weights must therefore be served mostly from on-chip caches, which is the data reuse that the roofline model~\cite{williams2009roofline} makes visible. The PCIe link is much slower than device memory (\val{Sp:micro:h2d_pinned_bandwidth:0}\,GB/s pinned against \val{Sp:micro:d2d_copy_bandwidth:0}\,GB/s), so data have to be reused on the device to pay for the transfer; with 3{,}136 bytes per image crossing the link once, the copy costs only \val{Sp:cuda:h2d:4096}\,ms per pass here~\cite{gregg2011data,harris2012transfers}. On the CPU side, the serial code also re-reads the weights for every image: this is \val{Sp:serial:wtraffic:4096}\,GB/s of weight traffic, below the \val{Sp:micro:cpu_triad_bandwidth:1}\,GB/s one thread achieves on the memory benchmark, and compiling with \texttt{-march=native} makes it only \val{X:serial:O3OverNative:4096}$\times$ faster than plain \texttt{-O3} (Table~\ref{tab:flags}). Both observations are consistent with a loop limited by data movement rather than arithmetic. With 24 threads the weight traffic is \val{Sp:omp:best:wtraffic:4096}\,GB/s, far above the memory bandwidth of \val{Sp:micro:cpu_triad_bandwidth:24}\,GB/s measured with 24 threads, so the weights are served from the caches there as well.

\subsection{Launch overhead and GPU utilisation}
An asynchronous kernel launch costs \val{Sp:micro:kernel_launch_async:0}\,$\mu$s and a launch followed by a synchronisation \val{Sp:micro:kernel_launch_sync:0}\,$\mu$s (Table~\ref{tab:micro}). For one image per batch the program spends \val{Sp:cuda:otherus:1}\,$\mu$s per batch on launches, API calls and synchronisation, but the kernels themselves take \val{Sp:cuda:kus:1}\,$\mu$s per batch. With this mapping a batch of one image has 256 useful threads out of 4{,}096 launched (16$\times$16 blocks), so most of the K40m's multiprocessors are idle and every thread executes a 784-step dependent chain. This is a latency-bound situation that a larger batch removes, and the kernel time per pass falls from \val{Sp:cuda:kernel:1}\,ms at batch 1 to \val{Sp:cuda:kernel:4096}\,ms at batch 4096.

\subsection{Which implementation is best}
Table~\ref{tab:best} summarises the measurements. For one image per batch and for batch 8, OpenMP on the multicore CPU is fastest. From batch 32 on, the pipelined hybrid is fastest, and its advantage over the second-best option is between \val{Sp:hyb:vscpu:32}$\times$ (batch 32) and about 2$\times$ at larger batches. \cudabase{} is never the fastest: its single-threaded preprocessing and its synchronous copies leave the GPU idle for part of every batch, which the hybrid removes. These rankings hold for both compiler settings (the \texttt{-O3} runs show the same order), but the batch size at which \cudabase{} (fixed 16$\times$16 block) overtakes the best OpenMP configuration moves from \val{Sp:cuda:crossover} (\texttt{-march=native}) to \val{SpO3:cuda:crossover} (plain \texttt{-O3}), because the preprocessing stage is faster with plain \texttt{-O3} (\val{SpO3:cuda:pre:4096}\,ms against \val{Sp:cuda:pre:4096}\,ms per pass) and the CPU baselines are slightly slower (Table~\ref{tab:flags}).

\begin{coltable}{Effect of the compiler setting on the Spartan node: time per pass (ms) with \texttt{-march=native} and with plain \texttt{-O3}, and their ratio.}{tab:flags}
\fitcol{\begin{tabular}{lrrrrrr}
\toprule
Implementation & \multicolumn{3}{c}{Batch 256} & \multicolumn{3}{c}{Batch 4096} \\
\cmidrule(lr){2-4}\cmidrule(lr){5-7}
 & native (ms) & \texttt{-O3} (ms) & ratio & native (ms) & \texttt{-O3} (ms) & ratio \\
\midrule
Serial & 1{,}235.2 & 1{,}296.2 & 1.05 & 1{,}235.4 & 1{,}295.6 & 1.05 \\
OpenMP, best thread count & 58.9 & 61.1 & 1.04 & 57.3 & 61.1 & 1.07 \\
CUDA-only (16$\times$16) & 55.1 & 51.4 & 0.93 & 46.7 & 42.7 & 0.91 \\
Hybrid-pipelined, best & 26.0 & 26.1 & 1.01 & 24.7 & 24.8 & 1.00 \\
\bottomrule
\end{tabular}
}
\end{coltable}
